\subsection{Investment manager selection}

These entries describe how an asset owner finds, evaluates, hires, contracts with and keeps or replaces managers. They draw on \hyperref[meth:manager-skill-evaluation]{evaluating skill}, \hyperref[meth:benchmark-selection-quality]{benchmarks} and \hyperref[meth:downside-appraisal-measures]{downside and drawdown measures}. Fee and filter numbers are checked in \texttt{trading\_notes.numerical.manager\_selection.manager\_selection\_checks}; all inputs are invented. US legal rules are stated as of 18 October 2026.

\begin{qfmethod}[Manager selection and due diligence: universe, quantitative and qualitative screens, investment and operational due diligence, Type I and Type II errors in hiring and firing, returns- and holdings-based style analysis, capture and drawdown in selection, monitoring]\label{meth:manager-due-diligence}
\qffield{Statement} \emph{Process.} A search starts from a broad universe and narrows it with qualitative factors (``the people and organizational structure, the firm's investment philosophy, the decision-making process, and the strength of its equity research'') and quantitative factors (comparisons with benchmarks and peers, measured style and valuation characteristics), looking at all times for ``consistency between a firm's stated philosophy and process and its actual practices''. A questionnaire covering organisation and people, philosophy and process, resources, performance and fees produces a short list, followed by interviews and on-site visits \citep[chap.~7, pp.~466--468]{managing-investment-portfolios-2007}. Bacon describes due diligence as the ``Four Ps'': people (ethics, experience, availability, past mistakes), process (articulated, documented, repeatable, independently risk-controlled), performance (ethically produced, GIPS-compliant, consistent with the process, measured by a measure chosen for a reason) and potential (capacity). He adds that a manager is chosen for its fit with the existing roster as well as in isolation, and is then monitored for underperformance, style drift, better alternatives and reputational damage \citep[3rd ed., 2023, chap.~8, pp.~489--491]{wide-bacon-attribution}. Operational matters (ownership, regulatory history, litigation, insurance, financial condition, key-person and backup procedures, references from lost clients) appear in the questionnaire's organisation section \citep[chap.~7, Example~7-15, pp.~469--470]{managing-investment-portfolios-2007}.

\emph{Hiring and firing as a statistical filter.} Bailey, Richards and Tierney take the null hypothesis that a manager has at best zero value added. A Type~I error keeps (or hires) a zero-value-added manager, rejecting a true null; a Type~II error fires (or does not hire) a positive-value-added one. A coarse filter commits more Type~I errors, a fine filter more Type~II, and ``Both Type I and Type II errors are expensive'', including transition costs \citep[chap.~12, pp.~778--781]{managing-investment-portfolios-2007}.

\emph{Style analysis.} Returns-based style analysis regresses a fund's returns on style or asset-class index returns with coefficients that ``totals to 100\%'' and are typically nonnegative, choosing the exposures that minimise the variance of the tracking difference. ``The objective \ldots is not to minimize either the average value of this difference or the sum of the squared differences.'' $R^2$ is attributed to style and the remainder to selection. In Sharpe's example an unconstrained regression gives economically impossible weights (for example $-69.5$ percent intermediate bonds) with $R^2$ 95.20, while quadratic programming gives about 70 percent value and 30 percent small stocks with $R^2$ 92.22 \citep[sections ``Style Analysis'' and ``Mutual Fund Styles'', Table~2]{wide-sharpe-style}. Holdings-based style analysis classifies each position by characteristics such as market capitalisation, valuation ratios and sector, and may capture style changes sooner. It is more data-intensive and needs classification choices. Returns-based analysis needs little data but may mis-characterise current style and depends on the indices chosen \citep[chap.~7, pp.~436 and~442--444, Exhibit~7-15]{managing-investment-portfolios-2007}.

\qffield{Practice} \emph{Capture and drawdown in selection.} Up- and down-capture show how a manager's returns participate in rising and falling markets, and drawdown and recovery show the depth and length of losses that a client must sit through (\hyperref[meth:downside-appraisal-measures]{definitions and limits}). In selection they test whether the pattern matches the stated process (a defensive manager should show low down-capture) and whether the drawdowns fit the owner's risk tolerance.

\emph{Style analysis, original} (36 invented months). A fund built as 55 percent large value, 30 percent small value and 15 percent cash plus a selection return gives constrained style weights of 1.4, 56.6, 0, 27.1 and 14.9 percent (large growth, large value, small growth, small value, cash), with $R^2$ 0.986. The weights minimise the variance of the tracking difference, so adding a constant to every fund return changes only the mean difference, not the weights. Minimising the mean squared difference instead would penalise the selection return's mean and give other weights. Recovering weights close to the construction does not by itself show that the right objective was minimised.

\emph{The filter, original.} Assume 40 percent of candidate managers add 2 percent a year, 40 percent add nothing and 20 percent lose 1 percent, the mix Bailey et al. imagine (p.~781), with 5 percent tracking error and normal returns. Keeping managers whose five-year $t$-statistic is positive retains 81.4 percent of good managers, half of zero-value managers and 32.7 percent of bad ones; the Type~I rate is 44 percent and Type~II 19 percent. Requiring $t>1.645$ cuts Type~I to 3.9 percent but keeps only 22.6 percent of good managers (Type~II 77 percent). The average value added of those kept rises from 0.99 to 1.55 points, but the sponsor must search far more often.

\qffield{Conditions} Past performance has weak predictive power: of 81 top-quartile managers in Russell's database in 1997, 16 remained there in 1998, 7 in 1999 and none in 2000--2002 \citep[chap.~7, p.~466]{managing-investment-portfolios-2007}. Style analysis assumes constant exposures over the window and well-chosen, mutually exclusive indices.

\qffield{Pitfalls} Hiring on recent returns. Accepting back-tested or model track records; Bacon notes that ``No one will publish back-test numbers if the result is negative'' (p.~491). Treating several flattering appraisal measures as corroboration: ``More than one appraisal measure is a red flag!'' (p.~490). Firing for short-run underperformance and paying transition costs to hire last period's winner.

\qffield{See also} \hyperref[meth:manager-philosophy-process-behaviour]{philosophy, process and behaviour}; \hyperref[meth:vehicle-contract-fee-analysis]{vehicles, contracts and fees}; \hyperref[ex:manager-selection-case]{selection case}.

\end{qfmethod}

\begin{qfmethod}[Evaluating a manager's investment philosophy and decision process, and behavioural effects on team decisions: the claimed inefficiency, research, portfolio construction, risk control, sell discipline, consistency, and mitigations for groupthink, overconfidence and herding]\label{meth:manager-philosophy-process-behaviour}
\qffield{Statement} The questionnaire's philosophy and process section asks for ``the firm's investment philosophy and the market inefficiency that it is trying to capture, along with any supporting evidence for this inefficiency''; the research process; the risk-management function and risk models; how adherence to the stated style and process is monitored; the stock-selection process, including unique information sources and how buy and sell decisions are made; and portfolio construction. The resources section asks how research is produced and fed into construction, which models are used, and how trading, turnover and costs are managed \citep[chap.~7, p.~468]{managing-investment-portfolios-2007}. A good long record by the same people following consistent disciplines is more informative than comparable statistics with manager turnover and shifts in orientation (p.~466).

\emph{Behaviour in teams.} Expertise does not remove bias; experts can be more overconfident than laymen when feedback is limited \citep[NBER working paper 9222 version, p.~15]{wide-barberis-thaler-survey}. As a team converges on a decision, ``public doubts about the wisdom of the planned move are gradually suppressed''. Klein's premortem, in which members write the history of an imagined failure before the decision is final, ``overcomes the groupthink'' \citep[chap.~24, ``The Premortem'']{thinking-fast-slow-2011}. Open discussion over-weights early, assertive speakers, so members should write their positions before discussion to ``decorrelate error'' (chap.~7, ``Exaggerated Emotional Coherence (Halo Effect)''). Effective committees perform System~2 functions, probing and challenging an individual's intuition \citep[p.~6]{behavioral-finance-second-generation-2019}. Career concerns push managers to herd: younger mutual fund managers are more likely to be fired after bad performance, so they ``hug benchmarks'' \citep[chap.~15, p.~503]{s7-ang}.

\qffield{Practice} Interpretation, an evaluation checklist:
\begin{itemize}
\item \emph{Philosophy}: a specific, plausible source of excess return (an inefficiency, a risk premium, a structural advantage) and a reason it persists; evidence predating the product launch.
\item \emph{Process}: research that produces forecasts with measurable skill; construction that converts them efficiently (transfer coefficient, \hyperref[res:fundamental-law-active-management]{fundamental law}); risk control independent of the portfolio managers; an explicit sell discipline; trading-cost measurement.
\item \emph{Consistency}: style analysis and attribution match the philosophy; turnover, concentration and active share match the claimed breadth and conviction.
\item \emph{Team}: decision rights clear; dissent recorded; independent written views before meetings; premortems for large positions; pre-set review triggers for losing positions (against the disposition effect); compensation that does not reward benchmark-hugging or short-term asset gathering.
\end{itemize}
Recommendation rule: hire only where the philosophy, the process evidence and the performance evidence point the same way.

\qffield{Conditions} Qualitative judgements are themselves exposed to halo effects: a persuasive presentation can colour the reading of the numbers, which argues for scoring each area independently before discussion.

\qffield{Pitfalls} Accepting a philosophy stated after the returns were known. Mistaking an articulate story for a repeatable process. Treating team stability as sufficient when the process itself has drifted.

\qffield{See also} \hyperref[def:behavioural-biases]{behavioural biases}; \hyperref[meth:manager-due-diligence]{due diligence}.

\end{qfmethod}

\begin{qfmethod}[Vehicles, contracts and fees: pooled vehicles against separately managed accounts, contract provisions, ad valorem and performance-based fees (asymmetric, symmetric or fulcrum, hurdles, high-water marks, caps, crystallisation, clawback), US limits on performance fees, and fee analysis]\label{meth:vehicle-contract-fee-analysis}
\qffield{Statement} \emph{Fee forms.} Ad valorem (assets under management) fees are a percentage of assets, often tiered. A simple performance-based fee is a base fee plus a share of performance over the benchmark, sometimes with a cap, which limits the incentive to take risk, and a high-water mark, which requires cumulative outperformance since the last performance fee was paid. Ad valorem fees are simple and predictable. Symmetric incentive fees ``that reduce as well as increase compensation'' may align interests but create revenue volatility for the manager. A one-sided fee ``conveys a call option to the investment manager'' \citep[chap.~7, p.~467]{managing-investment-portfolios-2007}; the call's value rises with volatility, which can encourage risk-taking (chap.~9, p.~596). Bacon contrasts asymmetric fees (a low base, a hurdle, a participation rate and perhaps a cap) with rare symmetric fees that rebate for underperformance, called fulcrum fees in the United States. He lists what a performance-fee agreement must specify: the return methodology; one benchmark; the hurdle (soft or hard); the high-water mark and its resetting; participation and rebate rates; the period; accrual and payment; the capital base; clawback; and arbitration. He recommends worked examples in the agreement and multi-year measurement periods \citep[3rd ed., 2023, chap.~3, pp.~59--63]{wide-bacon-attribution}. Crystallisation is the payment of an accrued performance fee; equalisation keeps investors who enter a pooled fund at different times from paying for performance they did not receive (pp.~61--64).

\emph{Three forms of performance-based fee} (Interpretation, grouping the sources above): (1) a base fee plus a share of return above a benchmark or hurdle, with no rebate below it, often capped (asymmetric); (2) a symmetric fulcrum fee that rises and falls with performance around the fee earned at benchmark performance; and (3) a share of absolute gains above a high-water mark, often with a hurdle, the hedge fund form \citep[chap.~17, pp.~581--582]{s7-ang}.

\emph{Contract types.} A segregated mandate is governed by an investment management agreement and its guidelines (benchmark, permitted assets, risk limits, reporting, fees, termination); a pooled vehicle is governed by its prospectus or offering and partnership documents, which the investor accepts rather than negotiates \citep[chap.~1, pp.~2--3]{managing-investment-portfolios-2007}. Interpretation: the agreement allows tailoring and quick termination but needs negotiation and monitoring; fund documents are standard and cheap to accept but fix liquidity terms and fees for all investors.

\emph{US law.} Section 205(a)(1) of the Investment Advisers Act bars advisory contracts that pay the adviser a share of capital gains or appreciation. Section 205(b)(2) exempts contracts with registered investment companies, and with others for assets over \$1 million, if the fee is based on asset value averaged over a specified period and increases and decreases ``proportionately with the investment performance'' relative to an appropriate index \citep[subsecs.~(a)--(c)]{usc-15-80b-5}. Rule 205-3 allows performance fees with ``qualified clients'' who meet dollar thresholds set by the most recent Commission order \citep[paras.~(a) and~(d)(1)]{cfr-17-275-205-3}. The order of 1~May 2026, effective 29~June 2026, raised the thresholds from \$1{,}100{,}000 to \$1{,}400{,}000 of assets under management and from \$2{,}200{,}000 to \$2{,}700{,}000 of net worth, without generally applying them to earlier contracts \citep[secs.~II--III, p.~23521]{qualified-client-order-2026}. On 6~October 2026 the Commission proposed broadening performance-based compensation for registered funds and adding accredited investors to the qualified-client definition. Comments are due 7~December 2026, so at the date of writing the proposal is not law \citep[summary and dates, p.~63676]{performance-compensation-proposal-2026}.

\qffield{Practice} \emph{Pooled against separate.} Interpretation. A pooled vehicle gives small investors diversification, low minimums and shared administration, but one portfolio for all, investor-level tax inefficiency and exposure to other investors' flows. A separately managed account gives ownership of the securities, transparency, with ``completely transparent'' trades \citep[chap.~15, p.~502]{s7-ang}, client-specific restrictions, tax-lot control and in-kind transition, at the cost of size minimums, separate custody and administration, and more monitoring. A registered fund's leverage and fee limits differ from a private fund's (\hyperref[def:asset-management-industry]{the industry}).

\emph{Fee analysis, original} (average assets 100 million a year; gross excess returns $+4$, $-3$, $+5$ percent).
\begin{itemize}
\item Ad valorem 0.50 percent: 0.50 a year, 1.50 points of assets in total.
\item Asymmetric fee, 0.25 percent base plus 20 percent of excess above a cumulative high-water mark, capped at 1.25 percent: 1.05, 0.25 and 0.65, a total of 1.95. Year~3 pays only on the 2 points above the mark; without the mark it would pay 1.25.
\item Fulcrum fee, 0.50 percent $\pm$10 percent of excess, bounded at 0.10 and 0.90: 0.90, 0.20 and 0.90, a total of 2.00.
\item Of the 6 points of cumulative excess, these take 25, 32.5 and 33.3 percent.
\item For a manager with no skill (excess normal with mean zero, standard deviation 5 percent), the asymmetric fee's expected value is 0.566 a year against a base of 0.25: the owner pays on average 0.32 for luck. The fulcrum fee's expectation is its fulcrum, 0.50.
\item On Ang's NAV path 100, 120, 110, 145 a 20 percent incentive with a high-water mark pays 4.0, nothing, then 5.0 on $145-120$ \citep[chap.~17, p.~582]{s7-ang}.
\end{itemize}

\qffield{Conditions} Fee outcomes depend on the order of returns, the high-water mark's reset terms, the measurement period and whether excess is measured gross or net of base fees; the example states its terms. The legal thresholds are US rules as at 18 October 2026.

\qffield{Pitfalls} Comparing fee schedules on one year's returns. Paying an asymmetric fee annually on a strategy whose skill can only be judged over years. Leaving the capital base, clawback or benchmark changes undefined, Bacon's common sources of dispute (pp.~61--63).

\qffield{See also} \hyperref[meth:manager-due-diligence]{due diligence}; \hyperref[def:asset-management-industry]{asset management industry}; \hyperref[ex:manager-selection-case]{selection case}.

\end{qfmethod}

\begin{qfexample}[Manager selection case: a fictional foundation choosing between two equity managers on skill evidence, uncertainty, due diligence, process, team behaviour, vehicle and contract]\label{ex:manager-selection-case}
\qffield{Statement} A fictional foundation (US-based, 300 million) is replacing a global equity manager with a 40 million mandate benchmarked to a broad float-adjusted index. It applies \hyperref[meth:manager-due-diligence]{due diligence}, \hyperref[meth:manager-skill-evaluation]{skill evaluation} and \hyperref[meth:vehicle-contract-fee-analysis]{fee analysis}. Information ratios are net of fees and annual; all data are invented \citep[method from][chap.~12, pp.~776--781]{managing-investment-portfolios-2007}.

\emph{Candidates.} Manager X: nine-year record, IR 0.62, tracking error 4 percent, team unchanged, returns-based style close to the benchmark with a modest value tilt, down-capture 0.92, maximum drawdown relative to the benchmark of $-7$ percent. Manager Y: three-year record, IR 1.10, tracking error 6 percent, launched by a portfolio manager who left a larger firm; style analysis shows a growth and momentum tilt that rose sharply in the last year; down-capture 1.08.

\qffield{Practice} \emph{Skill evidence and uncertainty.} The $t$-statistics are close: X $0.62\sqrt9=1.86$, Y $1.10\sqrt3=1.91$ (one-sided $p$ 0.031 and 0.028). Shrinking each IR toward zero with an assumed spread of true IRs of 0.30 and a sampling standard error of $1/\sqrt{\text{years}}$ gives posterior means of 0.277 for X and 0.234 for Y: Y's high ratio is less informative because its record is short. The ranking depends on that assumption. With a spread of 0.50 the posterior means are 0.429 for X and 0.471 for Y, and the two are equal at a spread of about 0.42. The quantitative evidence therefore does not separate the managers: X leads only if true IRs are believed to be tightly clustered near zero.

\emph{Due diligence and process.} X's attribution shows value from stock selection within sectors, which matches its stated research process; its risk control reports to the chief operating officer. Y's value came mostly from a growth factor tilt not described as its process, which is \hyperref[mech:biases-market-anomalies]{momentum} exposure obtainable more cheaply. Its sell decisions rest with one person. The foundation's committee wrote independent scores before discussing Y's presentation and ran a premortem, which surfaced key-person and style-drift risk.

\emph{Decision.} Hire X. The decision rests on the process, attribution and operational evidence, which favour X and do not depend on the prior, not on the posterior ranking, which does. Record Y as a candidate to monitor for two more years, when its longer record will make the ranking less sensitive to the prior.

\emph{Vehicle and contract.} 40 million justifies a separately managed account. It allows the foundation's exclusion list, holdings transparency and in-kind transition of part of the old portfolio. The contract is an investment management agreement with guidelines (benchmark, tracking-error range 2--6 percent, exclusions, cash limits), reporting, termination on 30 days' notice and a fee. The foundation prefers the 0.50 percent ad valorem fee to X's offered asymmetric 0.25 plus 20 percent alternative, because for a manager without skill the asymmetric fee is expected to cost more than the flat fee (0.566 against 0.50 in the fee example), and skill cannot be confirmed within a year.

\emph{Monitoring.} Quarterly reports; annual review against the IPS benchmark and the style analysis; a watch-list trigger on style drift, team departure or a three-year IR below zero, with review rather than automatic termination, to limit Type~II errors.

\qffield{Conditions} The prior spread of true IRs and the normal model are assumptions. A different prior changes the posterior means and, above a spread of about 0.42, the direction of the comparison, so the hiring decision must not depend on it.

\qffield{Pitfalls} Choosing Y on its ratio alone; choosing X without the process and attribution checks that make its ratio credible.

\qffield{See also} \hyperref[meth:manager-philosophy-process-behaviour]{process and behaviour}; \hyperref[meth:benchmark-selection-quality]{benchmarks}.

\end{qfexample}
